% Part of main.tex; compile main.tex.
% ---------------------------------------------------------------------------
\section{Related Work}
\label{sec:related}

In this section, we review asynchronous I/O frameworks for HPC, describe how
influential evaluations of these frameworks emulate the computation phase,
and summarize the existing evidence on interference between background I/O
and the application.

\subsection{Asynchronous I/O frameworks for HPC}

Asynchronous read and write operations have long been available through
POSIX AIO and non-blocking MPI-IO, and early studies examined overlap
strategies at the MPI-IO level~\cite{patrick2008overlap}. However, these
interfaces do not cover non-data operations such as opening and closing
files, and they require users to manage dependencies and byte offsets
manually~\cite{tang2022async}. For this reason, high-level libraries have
added their own mechanisms. Tang et al.\ introduced an HDF5 VOL connector
that queues operations, tracks their dependencies as a directed acyclic graph
(DAG), and executes them on a background thread managed by
Argobots~\cite{tang2019async,tang2022async}. Damaris dedicates one or more
cores per node to data management tasks, including asynchronous
I/O~\cite{dorier2016damaris}. In addition, staging-based and server-based
systems move data off the compute processes and persist it later.

\subsection{Sleep-based compute emulation}

Several influential evaluations of asynchronous I/O emulate the computation
phase with \texttt{sleep()}. We describe them below.

\paragraph{h5bench.} The h5bench suite~\cite{li2021h5bench} exposes an
\texttt{EMULATED\_COMPUTE\_TIME\_PER\_TIMESTEP} parameter. After each
iteration, the write and read kernels sleep for this duration to emulate the
application's computation. Because h5bench is widely used as a standard HDF5
benchmark, this choice propagates into many downstream studies.

\paragraph{HDF5 asynchronous VOL.} In the TPDS evaluation of the HDF5
asynchronous VOL connector~\cite{tang2022async}, all experiments write five
timesteps with a sleep between them, sized so that asynchronous I/O can fully
overlap with it. For the Nyx workload, all computation was replaced with
sleep, so that the same data structures could be written without running the
expensive simulation. The authors argue that the sleep durations, up to
60\,s on Cori and 20\,s on Summit, are conservative relative to real VPIC
compute phases.

\paragraph{Modeling asynchronous I/O.} Ravi et al.~\cite{ravi2023evaluating}
develop a model in which an asynchronous step costs
$\max(t_{comp}, t_{io} - t_{comp}) + t_{\text{transact overhead}}$. Their
VPIC-IO and BD-CATS-IO experiments use a 30-second sleep in place of the
computation, and the emulated compute length is tuned manually to completely
hide the asynchronous I/O time.

\paragraph{Proactive Data Containers.} Proactive Data Containers (PDC) move
data asynchronously through server processes that run on extra
cores~\cite{tang2018pdc}. In the CCGrid'18 evaluation of PDC, the
multi-timestep VPIC-IO runs ``added a sleep time, representing computations
between writing data,'' and the client observes only the last write plus
overhead. This yields up to a $5\times$ speedup over HDF5. The follow-up
tuning study~\cite{tang2019pdctune} configures Nyx and BD-CATS-IO with
``computation time (longer than I/O time)'' between timesteps, but does not
describe what that computation is. However, it does state that PDC reserves
one core per node for its servers.

\paragraph{ADIOS.} The BP5 engine of ADIOS2 can write buffered data in the
background after \texttt{EndStep}. It does so either naively or guided by
\texttt{EnterComputationBlock} and \texttt{ExitComputationBlock} hints
provided by the application~\cite{godoy2020adios2}. The ADIOS2 skeleton tool,
\texttt{adios\_iotest}, expresses the compute phase of a workflow as a
\texttt{sleep} command, or a \texttt{busy} loop, between I/O
steps~\cite{adios2_iotest}. The earlier Skel generator produces ADIOS
skeletal applications whose source contains ``the appropriate ADIOS I/O
calls, data generation, and timers'' and no computation at
all~\cite{logan2011skel}. It is important to note that the ADIOS developers
themselves report mixed results for asynchronous writes in real codes. They
observed a large gain for WarpX on Perlmutter, while on Frontier ``we don't
see improvement over synchronous I/O''~\cite{podhorszki2024adiosslides}. In
addition, the openPMD-api documentation warns that asynchronous writes ``can
compete with MPI for communication resources, impacting the \emph{compute}
performance of an application''~\cite{openpmd_adios2_docs}.

\paragraph{Deep-learning I/O.} DLIO~\cite{devarajan2021dlio} emulates the I/O
of scientific deep-learning training. Its per-step
\texttt{computation\_time} is profiled on a specific accelerator, and its
shipped configurations carry separate values for V100, A100, and H100 GPUs.
This time is replayed by calling \texttt{time.sleep()}~\cite{dlio_github}.
Because the DLIO reader workers prefetch the next batches while the main
process sleeps, the benchmark cannot observe host-side contention between the
input pipeline and the training step.

\paragraph{POSIX-level and checkpointing systems.} Systems that flush
checkpoints or burst-buffer contents in the background through POSIX are
rarely evaluated against a running computation. Gossman et al.\ evaluate
aggregation strategies for the asynchronous flush of VeloC by timing the
flush phase alone~\cite{gossman2021aggregated,nicolae2019veloc}. Harmonia,
an interference-aware burst-buffer scheduler, uses a synthetic benchmark
whose computation phases are ``artificial math kernels,'' and the content of
these kernels is not specified~\cite{kougkas2018harmonia}.

\subsection{Evidence on interference and overlap}

To the best of our knowledge, no study focuses specifically on sleep-based
emulation as a methodological problem. Instead, the relevant evidence is
distributed across several lines of work.

\paragraph{Caveats within sleep-based studies.} The HDF5 asynchronous VOL
paper notes that the background thread can affect application performance
if all cores are used by the application, and recommends sparing resources
for it~\cite{tang2022async}. It also explains that starting background
execution while the main thread is active can cause the two threads to
compete for shared resources, in the worst case serializing them. Ravi et
al.\ limit the scope of their model to the absence of node-level
oversubscription~\cite{ravi2023evaluating}. In addition, guidance from the
HDF Group recommends that OpenMP applications leave one core per process for
the background thread~\cite{hdfgroup2022sc}. In each case, the identified
risk is one that a sleeping main thread cannot exhibit.

\paragraph{Real-application evaluations.} Jain et al.\ evaluated HDF5
asynchronous I/O in Flash-X on Summit~\cite{jain2022flashx}. They note that
background threads compete for node resources and with communication, and
they observe a measurable overhead from asynchronous I/O. This overhead is
outweighed by the reduction in I/O time, which yields up to a $1.5\times$
overall speedup. This figure is far more modest than the I/O-time speedups of
roughly $4.5$ to $4.9\times$ reported with sleep-based
emulation~\cite{tang2022async}, although the two metrics are not directly
comparable.

\paragraph{Interference studies.} Tseng et al.~\cite{tseng2021demystifying}
study interference from background I/O directly. They argue that background
I/O competes with the application for CPU time, memory bandwidth, and
network bandwidth, and that multicore nodes make the problem worse because
many I/O operations compete concurrently. Without accounting for this,
asynchronous I/O may perform worse than blocking I/O. Their methodology
stresses CPU and memory bandwidth separately and is validated on HACC.
Tarraf et al.~\cite{tarraf2024behind} build on this observation. They note
that developers often assume overlap will be sufficient regardless of the
configuration, and they evaluate with real computation (WaComM++ and a
modified HACC-IO) rather than sleep.

\paragraph{Overlap benchmark methodology.} Benchmarks designed to measure
overlap generally use computation that occupies the CPU. The non-blocking
IMB-IO tests of the Intel MPI Benchmarks run a calibrated CPU-intensive
kernel concurrently with the I/O, and compute the overlap as
$(t_{pure} + t_{CPU} - t_{ovrl}) / \min(t_{pure}, t_{CPU})$~\cite{intel_imb_io}.
Similarly, COMB~\cite{lawry2002comb} and the MPI overlap suite of Denis and
Trahay~\cite{denis2016overlap} measure the overlap of communication and
computation against busy work. These designs assume that the compute side
can be slowed by the background activity, which is not possible with a
sleep.
